\section{Education and Outreach Plan}
Our Education and Outreach Plan aims to reach Chicago audiences, including K-12 students and educators, undergraduate and graduate students, and practitioners in the field of robotics and cybersecurity. 
The main goals are to offer educational resources, engage diverse communities, and inspire future careers in robotics, cybersecurity, and related fields.
% Educational robots can significantly enhance learning experiences by promoting active engagement, problem-solving, and collaboration among students. 
As part of our support for Chicago communities, the PI will focus on leveraging the unique opportunities and resources available in Chicago and Loyola University.   
% Since the PI is a faculty of Loyola University Chicago, the PI will develop robotics classes and provide guest lectures, and seminars. This will include several key initiatives:

\BfPara{Robotic Security Course Development} 
Loyola University Chicago has been designated as a Center of Academic Excellence (CAE) in Cyber Defense by the NSA/DHS in June 2020 and offers robust cybersecurity and computer science programs.
PI is an active member of the center and CAE community and will leverage these resources to design and develop a comprehensive robotic security course for university-level students. 
This course will cover a variety of topics, including basic robotics principles, programming, mechanical design and embodiments, AI integration, and emerging threats and countermeasures in robotics. 
% The material and outcomes of this project will be used to enrich the curriculum of undergraduate and graduate programs.
The material for this course will be competency-based and follows the DoD Cyber Workforce Framework and the NICE framework. Course materials and labs will be made available through the CLARK Cybersecurity Library.

\BfPara{Security Lectures and Panels} 
PI will host panels from industry and academia to discuss trends in the intersection of AI, robotics, and security. 
The PI is the lead organizer (PI) for the SecureAI program, an experiential cybersecurity training program for AI professionals, supported by the NSF.
Following these efforts, the PI
will also invite experienced professionals/researchers from the field of robotics and security to deliver guest lectures at Loyola University Chicago. 
These lectures will provide participants/students with insights into the latest advancements in robotics, real-world applications, and career opportunities. Topics will include emerging threats in robotics, security measures, and privacy concerns.

\BfPara{CS Fall Seminar Series} The PI hosts regular seminar series during the Fall semesters every year.
The material and efforts that originate from this project will be included in this series of seminars. The PI will ensure to feature topics related to robotic security and demo some projects to boost interest in the field.
%Several seminars will be organized to provide hands-on learning experiences for students. These events will focus on various aspects of robotics, such as building and programming robots, understanding AI and machine learning, and exploring the ethical implications of robotics technology. Workshops will include sessions on cybersecurity in robotics, addressing vulnerabilities, and implementing privacy protections to ensure students are aware of the critical aspects of security and privacy in their projects.

\BfPara{Undergraduate Research Participation} 
The PI Abuhamad is a co-PI of the CyberRamblers Scholarship for Service Program funded by the NSF. CyberRamblers scholars are mentored by the PI and encouraged to participate in research.
In addition to CyberRamblers, Loyola University Chicago has multiple initiatives for that can be leveraged for undergraduate research participation, such as \emph{LUROP program} and \emph{Provost, Mulcahy, and Carbon Scholarships/Fellowships}.
Students participating in research will access cutting-edge research facilities and opportunities to contribute to ongoing research in the field of robotics and cybersecurity.


%Research topics will include swarm robotics, soft robotics, and the integration of AI with a focus on security and privacy countermeasures.



\BfPara{Community Engagement and Outreach} 
(1) The PI will engage with the broader community through initiatives such as \emph{public lectures} and \emph{interactive workshops} hosted at Loyola University Chicago. These events will be designed to raise awareness and foster understanding of the critical role of robotics, AI, security, and privacy in  society. 
(2) In collaboration with the CS department, which has established strong partnerships with local schools and community organizations, the PI will launch a summer internship program for talented high school students. This program will provide participants with hands-on experience in robotic security research and development. 
(3) The PI plans to develop age-appropriate curriculum modules on robotics security and privacy for K-12 schools with a strong emphasis on security and privacy and leverage existing partnerships with Chicago Public Schools to integrate them into their STEM Summer Programs.


